\section{Packet normalization and polarization}

Let $T$ be the common linear transform used to assemble the source packets.  With the
project's inner-product convention, the four-phase identity has the form
\begin{equation}
 \langle T\beta,Tw\rangle
 =\frac14\sum_{j=0}^3 i^j\|T(\beta+i^jw)\|^2.          \label{eq:polarization}
\end{equation}
The key word is \emph{common}: the same linear map acts on all four inputs.

\begin{proposition}[Output normalization erases polarization]
Suppose all four vectors $T(\beta+i^jw)$ are nonzero.  If each output in
\eqref{eq:polarization} is
replaced by its own unit normalization, then the resulting right-hand side equals
zero.  It therefore cannot reproduce (6) in general.
\end{proposition}

\begin{proof}
Every normalized squared norm equals one, while
$\sum_{j=0}^3i^j=0$.  For a concrete counterexample take the one-dimensional map
$T=1$, $\beta=1$, and $w=2$.  The four raw squared norms are $9,5,1,5$ and \eqref{eq:polarization} equals
$2$, whereas their independently normalized replacements are $1,1,1,1$ and give
zero.
\end{proof}

Consequently, the normalized-row Bessel theorem from the model clock is not
automatically source-valid.  A fixed linear rescaling common to every packet may be
legal, but its inverse or compensating weight must remain visible in the physical
reassembly.  Output-dependent unit normalization cannot be hidden inside
\eqref{eq:full-source}.
